\section{Recovering the matrix obstruction}\label{sec:obstruction}

The preceding construction turns a small commutator into a small
one-sided gap. We now show that a spatial embedding close to the
identity would turn the original matrix pair into a nearby commuting
pair. Choosing the clock and shift matrices will make this impossible.

\subsection{The winding obstruction}
The matrices below are Voiculescu's almost commuting unitary pair
\cite{Voi1983}. We include the logarithmic winding argument, in the
form developed in \cite{EL1989,Exel1993}.

\begin{lemma}\label{lem:clock}
For $n\geq3$, let $U_n,V_n$ be the unitaries on $\mathbb C^n$ given by
\begin{equation}\label{eq:clock-shift}
 U_ne_j=\omega_n^je_j,\qquad
 V_ne_j=e_{j+1\bmod n},\qquad
 \omega_n=e^{2\pi i/n}\quad(0\leq j<n).
\end{equation}
Then $\nu_n:=\|U_nV_n-V_nU_n\|=|\omega_n-1|\to0$.
There is no infinite subsequence along which commuting unitaries
$\widehat U_n,\widehat V_n$ on $\mathbb C^n$ satisfy
\[
 \|\widehat U_n-U_n\|+\|\widehat V_n-V_n\|\longrightarrow0.
\]
\end{lemma}

\begin{proof}
The identity $U_nV_n=\omega_nV_nU_n$ proves the first assertion.
For any unitary matrices $U,V$ whose commutator $UVU^*V^*$ has no
eigenvalue $-1$, define
\[
 \kappa(U,V)=\frac{1}{2\pi i}\Tr\log(UVU^*V^*),
\]
using the principal logarithm and the unnormalized trace.
The determinant identity
\[
 \exp\bigl(\Tr\log(UVU^*V^*)\bigr)
       =\det(UVU^*V^*)=1
\]
shows that $\kappa(U,V)$ is an integer. It is continuous, and hence
constant, along any path on which the logarithm remains defined.
For the pair in \eqref{eq:clock-shift}, its value is $1$, since
$U_nV_nU_n^*V_n^*=\omega_n I$ and $0<2\pi/n<\pi$.
For a commuting pair its value is $0$.

Suppose the asserted approximants existed along a subsequence.
For all sufficiently large $n$ on that subsequence, write
$U_n^*\widehat U_n=e^{iA_n}$ and
$V_n^*\widehat V_n=e^{iB_n}$, with $A_n,B_n$ self-adjoint and
spectra in $(-\pi,\pi)$. The paths
\[
 U_n(s)=U_ne^{isA_n},\qquad
 V_n(s)=V_ne^{isB_n}\qquad(0\leq s\leq1)
\]
stay within the respective endpoint distances from $U_n,V_n$.
Indeed, for $|\theta|<\pi$ and $0\leq s\leq1$,
$|e^{is\theta}-1|\leq|e^{i\theta}-1|$.
A four-factor telescoping estimate gives, uniformly in $s$,
\[
 \|U_n(s)V_n(s)U_n(s)^*V_n(s)^*-I\|
 \leq \nu_n+2\|\widehat U_n-U_n\|
             +2\|\widehat V_n-V_n\|\longrightarrow0.
\]
The principal logarithm is therefore defined along the whole path
for all sufficiently large $n$. Its integer changes from $1$ to $0$,
a contradiction.
\end{proof}

\subsection{A finite-dimensional corner}
Compression of commuting operators to a prescribed subspace need not
preserve commutation. The next lemma avoids this difficulty by
replacing each support projection by a nearby projection in its
corresponding algebra and then using their common range.

\begin{lemma}\label{lem:corner}
For each $n$, let $E_{1n},E_{2n}$ be commuting projections on a complex
Hilbert space $H_n$, with $D_n=E_nH_n$ nonzero and finite-dimensional,
where $E_n=E_{1n}E_{2n}$. Let $\widetilde R_{1n},\widetilde R_{2n}$
be unitaries commuting with both projections, and write
$R_{in}=\widetilde R_{in}|_{D_n}$.
Suppose $\mathcal A_{1n},\mathcal A_{2n}\subset\B(H_n)$ are
commuting unital von Neumann algebras and there are
$T_{in}\in\mathcal A_{in}$ such that
\begin{equation}\label{eq:corner-hypothesis}
 \|T_{in}-E_{in}\widetilde R_{in}\|\longrightarrow0
 \quad(i=1,2).
\end{equation}
Then, for all sufficiently large $n$, there are commuting unitaries
$\widehat R_{1n},\widehat R_{2n}$ on $D_n$ with
\[
 \|\widehat R_{in}-R_{in}\|\longrightarrow0\quad(i=1,2).
\]
\end{lemma}

\begin{proof}
We first replace the intended projections by nearby projections in
the commuting algebras. Set
$\delta_{in}=\|T_{in}-E_{in}\widetilde R_{in}\|$.
Since $(E_{in}\widetilde R_{in})^*(E_{in}\widetilde R_{in})=E_{in}$,
\[
 d_{in}:=\|T_{in}^*T_{in}-E_{in}\|
           \leq \delta_{in}(2+\delta_{in})\longrightarrow0.
\]
For sufficiently large $n$, take the spectral projection
\[
 Q_{in}=\mathbf1_{(1/2,\infty)}(T_{in}^*T_{in})\in\mathcal A_{in}.
\]
The spectrum of $T_{in}^*T_{in}$ lies in
$[0,d_{in}]\cup[1-d_{in},1+d_{in}]$. This follows from the
resolvent bound for a self-adjoint operator within $d_{in}$ of a
projection. Functional calculus therefore gives
\[
 \|Q_{in}-T_{in}^*T_{in}\|\leq d_{in},
 \qquad \|Q_{in}-E_{in}\|\leq2d_{in}\longrightarrow0.
\]
Since the two algebras commute, $P_n=Q_{1n}Q_{2n}$ is a projection,
and
\[
 \|P_n-E_n\|
 \leq\|Q_{1n}-E_{1n}\|+\|Q_{2n}-E_{2n}\|\longrightarrow0.
\]

We identify this perturbed corner with $D_n$ by a small unitary.
The operator
\[
 S_n=P_nE_n+(I-P_n)(I-E_n)
     =I+(P_n-E_n)(2E_n-I)
\]
tends to the identity in norm and satisfies $S_nE_n=P_nS_n$.
It is invertible for large $n$. Its absolute value commutes with
$E_n$, so its unitary polar part
$J_n=S_n(S_n^*S_n)^{-1/2}$ satisfies
\begin{equation}\label{eq:corner-identification}
 J_nE_nJ_n^*=P_n,\qquad \|J_n-I\|\longrightarrow0.
\end{equation}
In particular, $P_nH_n$ has exactly the dimension of $D_n$.

The compressions $P_nT_{1n}P_n$ and $P_nT_{2n}P_n$ commute,
including across adjoints. This needs a check because $Q_{in}$ need
not commute with $T_{in}$. Suppress $n$ and put
$A=Q_1T_1Q_1\in\mathcal A_1$ and
$B=Q_2T_2Q_2\in\mathcal A_2$. The two compressions are $AQ_2$
and $Q_1B$. Since $AQ_1=Q_1A=A$, $BQ_2=Q_2B=B$, and every
element of $\mathcal A_1$ commutes with every element of
$\mathcal A_2$, both orders of their product equal $AB$.
Replacing $A$ or $B$ by its adjoint proves the other commutations.

Transfer these compressions to $D_n$ using \eqref{eq:corner-identification}:
\[
 C_{in}=\left.E_nJ_n^*T_{in}J_nE_n\right|_{D_n}.
\]
The pair still commutes across adjoints. Moreover,
\[
 \|C_{in}-R_{in}\|
 \leq\delta_{in}+2\|J_n-I\|\longrightarrow0,
\]
by comparing $T_{in}$ with $E_{in}\widetilde R_{in}$ before
conjugation. Thus $C_{in}$ is invertible for large $n$ and
$\|C_{in}^*C_{in}-I_{D_n}\|\to0$. Its unitary polar part
\[
 \widehat R_{in}=C_{in}(C_{in}^*C_{in})^{-1/2}
\]
remains close to $R_{in}$. The unital $C^*$-algebras generated by
$C_{1n}$ and $C_{2n}$ commute; continuous functional calculus puts
each polar part in its respective algebra. Hence the two polar
unitaries commute, as required.
\end{proof}

\subsection{Proof of the main theorem}
Apply the construction of Section~\ref{sec:near} to
$D=\mathbb C^n$ and the clock--shift pair of Lemma~\ref{lem:clock}.
Use subscripts $n$ for the resulting objects $L_n,H_n,a_n,b_n,f_n,W_n,
M_n,N_n$. All these Hilbert spaces are separable. By
Proposition~\ref{prop:near},
\begin{equation}\label{eq:vanishing-gap}
 \gamma(M_n,N_n)\leq C\nu_n\longrightarrow0.
\end{equation}

\begin{proof}[Proof of Theorem~\ref{thm:main}]
We show that no infinite subsequence admits unitaries $u_n$ with
\begin{equation}\label{eq:assumed-embedding}
 u_nM_nu_n^*\subset N_n,\qquad\|u_n-I\|\longrightarrow0.
\end{equation}
Suppose otherwise and work on that subsequence. Define, for $g\in G$,
\[
 Z_n(g)=W_n(\lambda(g)\otimes I_{L_n})W_n^*,\qquad
 \Sigma_n(g)=u_nZ_n(g)u_n^*(\rho(g)\otimes I_{L_n}).
\]
The first factor in the expression for $\Sigma_n(g)$ lies in $N_n$,
by \eqref{eq:assumed-embedding}. Thus the two representations being
multiplied have commuting ranges, and $\Sigma_n$ is a unitary
representation of $G$. Proposition~\ref{prop:near} gives the uniform
estimate
\begin{equation}\label{eq:sigma-bound}
 \sup_{g\in G}
 \bigl\|\Sigma_n(g)-\lambda(g)\rho(g)\otimes f_n(g)\bigr\|
 \leq\beta_n,
 \qquad \beta_n=C\nu_n+2\|u_n-I\|\longrightarrow0.
\end{equation}
Let $\mathcal A_{1n}$ and $\mathcal A_{2n}$ be the von Neumann
algebras generated by the restrictions of $\Sigma_n$ to $F\times\{e\}$
and $\{e\}\times F$, respectively. These algebras commute.

We will apply Lemma~\ref{lem:corner} by taking weak limits of the
generators in each factor. Under $S=K\otimes K$, write
\[
 H_n=K\otimes K\otimes K\otimes K\otimes\mathbb C^n.
\]
Let $E_{1n}$ use $p_e$ in the first and third slots, and let $E_{2n}$
use $p_e$ in the second and fourth slots, with identities in all
other slots. These projections commute, and
$E_n=E_{1n}E_{2n}$ has range naturally identified with $\mathbb C^n$.
Let $\widetilde R_{1n}$ and $\widetilde R_{2n}$ act as $U_n,V_n$,
respectively, on the final slot and as the identity elsewhere.

Set $g_{1j}=(s_j,e)$ and $g_{2j}=(e,s_j)$. On the $i$-th factor
of $S$, the operator $\lambda(g_{ij})\rho(g_{ij})$ sends
$\delta_x$ to $\delta_{s_jxs_j^{-1}}$. For each fixed nonidentity
word $x$, this conjugate leaves every finite subset of $F$ as
$j\to\infty$: for large $j$, the letter $s_j$ occurs neither in
$x$ nor in the finitely many words under consideration. Hence the
conjugation operators converge weakly to $p_e$ in that slot.
Combining this observation with \eqref{eq:matrix-limits} gives,
for each fixed $n$ and $i$,
\begin{equation}\label{eq:joint-limit}
 \lambda(g_{ij})\rho(g_{ij})\otimes f_n(g_{ij})
 \xrightarrow[j\to\infty]{\WOT}tE_{in}\widetilde R_{in}.
\end{equation}
Here tensor products of bounded weakly convergent sequences on
separate slots converge weakly, as follows first on elementary
tensors and then by density. The finite spectral decompositions
used in \eqref{eq:matrix-limits} justify the matrix-valued factor.

For each $i,n$, choose a weak operator cluster point $B_{in}$ of
$\Sigma_n(g_{ij})$ as $j\to\infty$. Such a point exists by weak
operator compactness of the closed unit ball of $\B(H_n)$, and
belongs to $\mathcal A_{in}$. Norm balls are weak operator closed,
so \eqref{eq:sigma-bound} and \eqref{eq:joint-limit} imply
\[
 \|B_{in}-tE_{in}\widetilde R_{in}\|\leq\beta_n.
\]
Thus $T_{in}=t^{-1}B_{in}\in\mathcal A_{in}$ satisfies
\eqref{eq:corner-hypothesis}. No product of weak limits is used:
commutation comes from membership in the two commuting algebras.
Lemma~\ref{lem:corner} now produces commuting unitary matrices on
$\mathbb C^n$ approaching $U_n,V_n$ in norm. This contradicts
Lemma~\ref{lem:clock}, proving that \eqref{eq:assumed-embedding}
cannot hold on any infinite subsequence.

It follows that some $\epsilon_0>0$ bounds
$\|u-I_{H_n}\|$ from below for every implementing unitary and all
sufficiently large $n$. Otherwise one could choose increasing
$n_j$ and implementing unitaries with distance less than $1/j$,
contrary to what we have just proved. Together with
\eqref{eq:vanishing-gap}, this is Theorem~\ref{thm:main}.
\end{proof}
